% Created 2026-02-22 Sun 22:25
% Intended LaTeX compiler: pdflatex
\documentclass[11pt]{article}
\usepackage[utf8]{inputenc}
\usepackage[T1]{fontenc}
\usepackage{graphicx}
\usepackage{longtable}
\usepackage{wrapfig}
\usepackage{rotating}
\usepackage[normalem]{ulem}
\usepackage{amsmath}
\usepackage{amssymb}
\usepackage{capt-of}
\usepackage{hyperref}
\author{Prologos Project — from Zachary Larson}
\date{2026-02-22}
\title{Extended Spec Language Design\\\medskip
\large Toward Editor-Assisted Interactive Development}
\hypersetup{
 pdfauthor={Prologos Project — from Zachary Larson},
 pdftitle={Extended Spec Language Design},
 pdfkeywords={},
 pdfsubject={},
 pdfcreator={Emacs 30.2 (Org mode 9.7.39)}, 
 pdflang={English}}
\begin{document}

\maketitle
\tableofcontents

\section{Executive Summary}
\label{sec:org9c2dd4c}

This document proposes extending Prologos's surface language with three
coordinated features that make dependent types, properties, and higher-kinded
abstractions approachable without sacrificing formal rigor:

\begin{enumerate}
\item \textbf{Extended \texttt{spec}} with keyword metadata (\texttt{:implicits}, \texttt{:where}, \texttt{:doc},
\texttt{:examples}, \texttt{:properties}, contracts, refinements)
\item \textbf{\texttt{property} keyword} for reusable, composable proposition groups (analogous
to \texttt{bundle} for traits), with \texttt{:includes} composition and \texttt{:laws} on traits
\item \textbf{\texttt{functor} keyword} for named type abstractions with category-theoretic
metadata (\texttt{:unfolds}, \texttt{:compose}, \texttt{:identity}, \texttt{:laws}) --- eliminating
raw Pi/Sigma types from the surface language
\item \textbf{\texttt{??} typed holes} for interactive, hole-driven development
\end{enumerate}

The design philosophy: \emph{dependent types are the implementation substrate; the
surface language speaks in domain terms}. Pi types, Sigma types, universe
levels, and multiplicity annotations exist in \texttt{:unfolds} blocks, in the
elaborator, and in the type checker. The primary programmer interface is
structured metadata expressed through keywords.

The design is informed by a survey of Clojure Spec, Malli, PropEr, Agda, Idris,
Lean 4, Haskell, F*, Koka, and Racket contracts. Property-based testing and
interactive theorem proving are explicit design \emph{targets} (we design syntax to
support them) but not implementation targets (we do not build the engines yet).
\subsection{Keyword Symmetry Table}
\label{sec:orgcbaef95}

\begin{center}
\begin{tabular}{llll}
Abstraction layer & Keyword & What it names & Analogous to\\
\hline
Function interface & \texttt{spec} & Function type + metadata & ML type signatures\\
Method requirement & \texttt{trait} & A single method signature & Haskell typeclass method\\
Requirement group & \texttt{bundle} & Conjunction of traits & Haskell superclass constraint\\
Proposition group & \texttt{property} & Conjunction of \texttt{:holds} clauses & QuickCheck property suites\\
Type abstraction & \texttt{functor} & Parameterized type + algebraic structure & CT functor / Lean \texttt{abbrev}, but richer\\
\end{tabular}
\end{center}
\section{Part I: Research Survey}
\label{sec:org527adb9}

\subsection{1. Clojure Spec \& Malli}
\label{sec:orga562454}

\subsubsection{Clojure Spec}
\label{sec:org99aa5bc}

Clojure Spec (\texttt{clojure.spec.alpha}) is \emph{predicate-first}: every spec is
fundamentally a predicate function. The key insight for Prologos is the three
facets of \texttt{s/fdef}:

\begin{verbatim}
(s/fdef ranged-rand
  :args (s/cat :start int? :end int?)
  :ret  int?
  :fn   (fn [{:keys [args ret]}]
          (and (>= ret (:start args))
               (<  ret (:end args)))))
\end{verbatim}

\begin{itemize}
\item \textbf{\texttt{:args}}: regex spec for function arguments (\texttt{s/cat}, \texttt{s/alt})
\item \textbf{\texttt{:ret}}: spec for the return value
\item \textbf{\texttt{:fn}}: predicate relating \texttt{\{:args conformed-args :ret conformed-ret\}}
\end{itemize}

The \texttt{:fn} relation is the most interesting for Prologos --- it expresses
\emph{properties} that must hold between inputs and outputs. In a dependently-typed
language, this corresponds directly to a dependent Pi type.

\textbf{Generative testing}: \texttt{stest/check} generates random inputs from \texttt{:args},
invokes the function, and checks both \texttt{:ret} and \texttt{:fn}. Specs \emph{are} generators
--- no separate test-writing step required. This is QuickCheck-style property
testing with zero user-written generators.

\textbf{Instrument vs check}: \texttt{instrument} validates \texttt{:args} at call sites
(development-time contract enforcement). \texttt{check} validates \texttt{:ret} and \texttt{:fn} via
generative testing (test-time property verification). This separation is
important: \emph{contracts at dev time, properties at test time}.

\textbf{Composition}: \texttt{s/keys} for entity maps, \texttt{s/cat} for sequential concatenation,
\texttt{s/alt} for alternatives, \texttt{s/and} for conjunctive refinement. Specs compose
algebraically.
\subsubsection{Malli}
\label{sec:org8d95780}

Malli is \emph{schema-first} (data-driven): schemas are plain data structures, not
opaque predicate objects. This enables schema transformation, generation,
serialization.

\begin{verbatim}
(def MyFunction
  [:=> [:cat :int :int] :int])

;; Schema-as-data enables programmatic manipulation:
(m/children MyFunction) ;; => [[:cat :int :int] :int]
\end{verbatim}

\textbf{Key lesson}: In a homoiconic language like Prologos, spec metadata should be
data --- quotable, inspectable, transformable. This is naturally satisfied by
using maps.
\subsubsection{Insight for Prologos}
\label{sec:orgd66419d}

In a dependently-typed language, types already subsume specs. Clojure's \texttt{:fn}
relation corresponds to a dependent Pi type:

\begin{verbatim}
;; Clojure spec's :fn relation for ranged-rand:
;; (fn [{:keys [args ret]}] (and (>= ret start) (< ret end)))

;; Is expressible as the dependent type:
;; <(start : Int) -> (end : Int) -> {v : Int | v >= start, v < end}>
\end{verbatim}

Runtime property testing remains valuable even when types encode the invariant:
the type checker catches \emph{logical} errors at compile time; property testing
catches \emph{implementation} bugs.
\subsection{2. PropEr (Erlang Property-Based Testing)}
\label{sec:org5592240}

PropEr is notable for its \emph{type-driven generator synthesis}. If a function has a
Dialyzer type spec, PropEr constructs generators automatically:

\begin{verbatim}
-spec foo(integer(), string()) -> boolean().
%% PropEr generates random (integer, string) pairs and checks boolean result
\end{verbatim}

\textbf{Compositional shrinking}: When custom generators are built from other
generators, shrinking behavior inherits automatically. Data structures
containing other data "tend to empty themselves."

\textbf{Property patterns}: \texttt{?FORALL(Xs, Generator, Property)} is the canonical form.
\texttt{?IMPLIES} for conditional properties. \texttt{aggregate/collect} for coverage.

\textbf{Insight for Prologos}: Our types are richer than Dialyzer specs. A \texttt{spec} with
dependent types gives enough information to auto-generate random inputs. The
generator infrastructure should be type-directed: \texttt{Gen : Type -> Generator} as a
trait, with automatic derivation for algebraic data types.
\subsection{3. Agda/Idris Interactive Development}
\label{sec:org558e5a3}

\subsubsection{Agda}
\label{sec:org6f69a3e}

Agda's hole-driven development is the gold standard for interactive type-driven
programming:

\begin{itemize}
\item \textbf{Holes}: \texttt{?} in code creates a \emph{hole} (placeholder for incomplete code). The
type checker reports the expected type and available local bindings.

\item \textbf{Commands}:
\begin{center}
\begin{tabular}{lll}
Key & Action & Description\\
\hline
\texttt{C-c C-r} & Refine & Fill hole with expression + new sub-holes\\
\texttt{C-c C-c} & Case & Split on a variable (pattern match)\\
\texttt{C-c C-a} & Auto & Proof search (try constructors)\\
\texttt{C-c C-,} & Goal & Show expected type and context\\
\end{tabular}
\end{center}

\item \textbf{Auto proof search}: Completely rewritten in Agda 2.7.0. Searches for type
inhabitants by trying constructors and projections. Works for "small enough
problems of almost any kind." Users provide hints via constants or scope.
\end{itemize}
\subsubsection{Idris 2}
\label{sec:org099d2d2}

Idris's interactive editing provides similar capabilities with a different
flavor:

\begin{itemize}
\item \textbf{Holes}: \texttt{?name} syntax. The REPL shows expected type and context.

\item \textbf{Commands}:
\begin{center}
\begin{tabular}{ll}
Command & Description\\
\hline
\texttt{:ac} & Add clause template for a function\\
\texttt{:cs} & Case split on variable (removes impossible cases)\\
\texttt{:am} & Add missing clauses\\
\texttt{:ps} & Proof search\\
\texttt{:mw} & Make with clause\\
\end{tabular}
\end{center}

\item \textbf{Case splitting} (\texttt{:cs}): Removes \emph{impossible} cases due to unification ---
this is type-directed case analysis. When a constructor can't unify with the
expected type, the case is silently omitted.
\end{itemize}
\subsubsection{Insight for Prologos}
\label{sec:org2e45f04}

The \texttt{??} syntax naturally maps to our existing \texttt{expr-hole} infrastructure
(\texttt{syntax.rkt} line 688). The type checker already handles \texttt{expr-hole} --- it
accepts any type. The missing piece is the \emph{interactive protocol}: the editor
asks "what type is expected here, what bindings are in scope?" and the type
checker reports back.

Key distinction: \texttt{\_} means "infer this for me automatically" (silent, automatic).
\texttt{??} means "I don't know what goes here, please help me" (interactive,
informational). The former is for the machine; the latter is for the human.
\subsection{4. Dependent Type Ergonomics}
\label{sec:orgd0ed6d4}

\subsubsection{How Languages Present Pi Types}
\label{sec:orgbc0ab8a}

\begin{center}
\begin{tabular}{llll}
Language & Pi syntax & Implicit & Instance\\
\hline
Agda & \texttt{(x : A) -> B} & \texttt{\{x : A\}} & \texttt{\{\{x : A\}\}}\\
Idris 2 & \texttt{(x : A) -> B} & \texttt{\{x : A\}} & \texttt{\{auto x : A\}}\\
Lean 4 & \texttt{(x : A) -> B} & \texttt{\{x : A\}} & \texttt{[x : A]}\\
Prologos & \texttt{<(x : A) -> B>} & \texttt{\{A : Type\}} & where clause\\
\end{tabular}
\end{center}

\textbf{The ergonomic tension}: Pi types with many parameters become hard to read.
All dependently-typed languages solve this the same way: separate the
\emph{signature} from the \emph{definition}. Prologos's \texttt{spec\textasciitilde{}/\textasciitilde{}defn} split already
achieves this. The extended spec pushes further by structuring the metadata
around the signature.
\subsubsection{Lean's \texttt{abbrev} vs \texttt{def} Distinction}
\label{sec:org8db01bb}

Lean 4 distinguishes transparent abbreviations from opaque definitions:

\begin{itemize}
\item \texttt{abbrev} is tagged \texttt{@[reducible]} --- unfolded during elaboration, instance
synthesis, and definitional equality checks. \texttt{abbrev MyMonad := ReaderT Env IO}
means \texttt{Monad MyMonad} resolves automatically via \texttt{Monad (ReaderT ...)}.
\item \texttt{def} (semireducible) is \emph{not} unfolded by instance synthesis. You need
explicit instances for opaque type definitions.
\end{itemize}

This distinction is critical for type aliases that should "carry through" trait
instances vs abstract types that deliberately hide their representation.
\subsubsection{Type-Level Functions in Dependently-Typed Languages}
\label{sec:orgad9ce56}

In Agda, Idris, and Lean, type synonyms are ordinary functions returning \texttt{Type}.
No special mechanism needed. The tradeoff: full type-level computation makes
inference harder and error messages may show fully-reduced forms rather than the
user's abbreviation name.

Haskell's \texttt{type} synonyms are transparent but \emph{cannot be partially applied} ---
a fundamental restriction. Prologos should avoid this limitation.
\subsubsection{Insight for Prologos}
\label{sec:orgd1994f9}

Our existing implicit argument inference (Sprint 3) and trait resolution (Phases
A--E) provide the core mechanisms. The ergonomic challenge is \emph{presentation} ---
how to show complex Pi types without overwhelming users. The extended \texttt{spec}
addresses function-level presentation; the \texttt{functor} keyword (Part VII)
addresses type-level presentation.
\subsection{5. Refinement Types (Liquid Haskell, F*)}
\label{sec:org3acded3}

\subsubsection{Liquid Haskell}
\label{sec:org38eba1d}

Refinement types refine base types with logical predicates:

\begin{verbatim}
{v : T | predicate}
\end{verbatim}

For functions, the return type can reference the input:

\begin{verbatim}
incr :: x:Int -> {v:Int | v = x + 1}
\end{verbatim}

Type checking reduces to SMT validity queries (static verification, not runtime
testing). Refinement reflection allows the SMT solver to reason about recursive
function implementations.
\subsubsection{F*}
\label{sec:org39f13fd}

Combines ML-like programming with dependent types, monadic effects, and
refinement types. Uses Z3 for automated verification. Programs extract to
OCaml, F\#, or C.

F* provides fine-grained unfolding control: \texttt{unfold} (always reduce),
\texttt{unfold\_for\_unification\_and\_vcgen} (reduce only during unification and VC
generation), or opaque (default). This three-tier system gives precise control
over where abbreviations are transparent.
\subsubsection{Insight for Prologos}
\label{sec:orgeb03883}

Prologos already has full dependent types, which \emph{subsume} refinement types:
\texttt{\{v : Int | v > 0\}} in Liquid Haskell is expressible as
\texttt{<(v : Int) * (GT v 0)>} (Sigma type pairing value with proof) in Prologos.

The ergonomic question: should we offer a lighter-weight refinement syntax
(\texttt{:refines} key) that \emph{compiles} to Sigma types? Yes --- for pragmatism. Users
write \texttt{:refines (fn [r] [>= r 0])} and the compiler generates the Sigma type.
\subsection{6. Contract Systems \& Spec Language Design Patterns}
\label{sec:org5fdfeba}

\subsubsection{Racket Contracts}
\label{sec:orga7e7e90}

\texttt{->i} for dependent contracts with \texttt{\#:pre} and \texttt{\#:post} conditions. \texttt{define/contract}
establishes a contract boundary. Blame tracking identifies which party violated
the contract.
\subsubsection{Composability Patterns}
\label{sec:org94d7b6a}

\begin{itemize}
\item The Specification Pattern (Fowler) composes predicates via AND/OR/NOT combinators.
\item The "schema explosion" problem occurs when specs grow combinatorially. Solution:
keep specs \emph{open} (new keys without modifying existing specs) and \emph{compositional}
(specs combine via simple operations). Clojure handles this via \texttt{s/merge} and \texttt{s/and}.
\end{itemize}
\subsubsection{Koka Effect Aliases}
\label{sec:org57ffbc6}

Koka's \texttt{alias pure = <div,exn>} abbreviates a row of effects. This demonstrates
the general principle: named groupings of constraints improve readability.
Prologos's \texttt{bundle} already does this for trait constraints; \texttt{property} will do
this for propositions; \texttt{functor} will do this for type structure.
\subsubsection{Insight for Prologos}
\label{sec:orgb8a7d98}

The map argument to \texttt{spec} is inherently open and composable. New keys can be added
in future versions without breaking existing code. Unrecognized keys are stored
but not acted upon (\emph{Postel principle}: liberal in what you accept, conservative
in what you emit). Spec metadata should be first-class data (homoiconicity
principle) --- quotable, inspectable, transformable.
\subsection{7. Type Abbreviation Mechanisms (Cross-Language Survey)}
\label{sec:org86b98e3}

\subsubsection{Transparent Aliases}
\label{sec:org2d0c29d}

Haskell \texttt{type}, Rust \texttt{type}, Koka \texttt{alias}. Fully expanded before type checking;
zero impact on inference. Haskell's limitation: type synonyms cannot be partially
applied. Rust and Koka do not share this restriction.
\subsubsection{Opaque / Controlled-Reducibility Aliases}
\label{sec:org68ad41c}

Lean 4's \texttt{abbrev} (reducible, instances carry through) vs \texttt{def} (semireducible,
opaque to instance search). Scala 3's \texttt{opaque type} (transparent inside defining
scope, opaque outside). Haskell's \texttt{newtype} (distinct type, zero-cost at runtime).
\subsubsection{Module-Level Abstraction}
\label{sec:orgf845975}

OCaml modules control type visibility through signatures: same \texttt{type t = int}
can be exposed transparently or abstractly depending on the \texttt{.mli} file. 1ML
(Rossberg) unifies core and module language so type abbreviations are ordinary
first-class functions.
\subsubsection{Insight for Prologos}
\label{sec:org1f5afa5}

The \texttt{functor} keyword (Part VII) takes the Lean \texttt{abbrev} approach --- transparent
by default, with the added richness of category-theoretic metadata. Unlike bare
abbreviations, a \texttt{functor} declares what the type \emph{means} through its operations,
not just what it expands to.
\section{Part II: Syntax Design}
\label{sec:orga409653}

\subsection{Current Spec Syntax}
\label{sec:orgf7c6e8f}

\begin{verbatim}
;; Basic:
spec add Nat Nat -> Nat

;; With implicit binders:
spec id {A : Type} A -> A

;; With trait constraints (where clause):
spec sort {A : Type} [List A] -> [List A] where (Ord A)

;; With inline constraints:
spec elem {A} [Eq A] A [List A] -> Bool

;; Multi-arity:
spec zip
  | {A B : Type} [List A] [List B] -> [List [Pair A B]]
  | {A B C : Type} [A -> B -> C] [List A] [List B] -> [List C]
\end{verbatim}

The \texttt{spec-entry} struct (macros.rkt line 203) stores: \texttt{type-datums},
\texttt{docstring}, \texttt{multi?}, \texttt{srcloc}, \texttt{where-constraints}, \texttt{implicit-binders},
\texttt{rest-type}.
\subsection{Extended Spec Syntax}
\label{sec:orga45539c}

The extension uses Prologos's existing implicit map syntax. Keyword-headed
children after the type signature are collected into a metadata map:

\begin{verbatim}
spec sort [List A] -> [List A]
  :implicits {A : Type}
  :where (Ord A)
  :doc "Sorts a list in ascending order using merge sort"
  :examples
    - [sort '[3N 1N 2N]] => '[1N 2N 3N]
    - [sort '[]] => '[]
  :properties (sortable-laws A)
  :see-also [reverse filter]
\end{verbatim}
\subsubsection{The \texttt{:implicits} Key}
\label{sec:orge7189f4}

The \texttt{:implicits} key lifts implicit binders out of the type signature into
metadata, freeing the signature to express pure function shape:

\begin{verbatim}
;; Current (inline implicits):
spec gmap {A B : Type} {C : Type -> Type} (Seqable C) -> (Buildable C) -> [A -> B] -> [C A] -> [C B]

;; With :implicits + :where:
spec gmap [A -> B] -> [C A] -> [C B]
  :implicits {A B : Type} {C : Type -> Type}
  :where (Seqable C) (Buildable C)
\end{verbatim}

All three forms coexist:
\begin{enumerate}
\item Inline \texttt{\{A : Type\}} --- for quick one-liners
\item \texttt{:implicits} key --- for clean multi-line specs
\item Omitted entirely --- when inference handles it (already works for kind \texttt{Type})
\end{enumerate}

\textbf{Implementation}: In \texttt{process-spec}, extract \texttt{:implicits} from the metadata map
and merge with any inline \texttt{\{...\}} binders. The \texttt{extract-implicit-binders}
function already returns \texttt{(values implicit-binders remaining-tokens)} --- union
the two sources. Duplicate binders (same name in both) emit a warning and
deduplicate.
\subsubsection{Multi-Arity Specs with Metadata}
\label{sec:orgbdade3f}

\textbf{Decision}: metadata is function-level (shared across all arities). Each \texttt{|}
branch defines its own type; the metadata applies to the function as a whole.

\begin{verbatim}
spec greet
  | Nat -> String
  | String Nat -> String
  :doc "Greet someone"
  :examples
    - [greet 42N] => "hello 42"
    - [greet "world" 3N] => "hello world world world"
  :properties
    - :name "greet-nat-positive"
      :forall {n : Nat}
      :holds [str-length [greet n] > 0N]
\end{verbatim}

Rationale: properties have \texttt{:forall} binders whose types determine which
branch they exercise. Examples are concrete calls whose argument types
determine dispatch. Per-branch metadata is theoretically more precise but
practically redundant --- the types already carry the dispatch information.
\subsubsection{Sexp Mode Equivalent}
\label{sec:org0e1cc49}

\begin{verbatim}
(spec sort [List A] -> [List A]
  (:implicits {A : Type})
  (:where (Ord A))
  (:doc "Sorts a list in ascending order using merge sort")
  (:examples (([sort '[3N 1N 2N]] => [1N 2N 3N])
              ([sort '[]] => [])))
  (:properties (sortable-laws A)))
\end{verbatim}
\subsubsection{How Parsing Works}
\label{sec:org05fda01}

\begin{enumerate}
\item The WS reader produces the outer form with keyword-headed children as nested
s-expressions.

\item \texttt{rewrite-implicit-map} (macros.rkt line 2452) is extended to recognize \texttt{spec}
as a trigger head alongside \texttt{def\textasciitilde{}/\textasciitilde{}defn}. The function detects keyword-headed
children \emph{after} the type signature and wraps them into a \texttt{\$brace-params}
(map literal) node.

\item \texttt{process-spec} (macros.rkt line 1314) is extended to look for a trailing map
after extracting the type signature. Recognized keys are extracted; the entire
map is stored in a new \texttt{metadata} field on \texttt{spec-entry}.

\item Backward compatibility: specs without keyword children parse identically to
today. The existing positional \texttt{where} keyword still works; \texttt{:where} in the
map augments it (union of both constraint sets).
\end{enumerate}
\subsection{Supported Keys}
\label{sec:org6df9fa7}

\subsubsection{Phase 1 (Near-term: syntax + storage)}
\label{sec:org3643dee}

\begin{center}
\begin{tabular}{lll}
Key & Type & Semantics\\
\hline
\texttt{:implicits} & \texttt{(listof brace-param-group)} & Implicit type binders (moved from inline)\\
\texttt{:where} & \texttt{(listof constraint)} & Trait constraints (migrated from positional)\\
\texttt{:doc} & \texttt{string} & Documentation string\\
\texttt{:examples} & \texttt{(listof (datum => datum))} & Input/output pairs for auto-test generation\\
\texttt{:see-also} & \texttt{(listof symbol)} & Cross-references to related specs\\
\texttt{:since} & \texttt{string} & Version/date introduced\\
\texttt{:deprecated} & \texttt{string} or \texttt{\#t} & Deprecation notice\\
\end{tabular}
\end{center}
\subsubsection{Phase 2 (Medium-term: property checking)}
\label{sec:org0add3e9}

\begin{center}
\begin{tabular}{lll}
Key & Type & Semantics\\
\hline
\texttt{:properties} & \texttt{(listof property-ref)} & Property references or inline property clauses\\
\texttt{:pre} & \texttt{datum} (args -> Bool) & Precondition (contract on arguments)\\
\texttt{:post} & \texttt{datum} (args ret -> Bool) & Postcondition (contract on return)\\
\texttt{:invariant} & \texttt{datum} (\{args ret\} -> Bool) & Relation between args and return (Clojure \texttt{:fn})\\
\end{tabular}
\end{center}
\subsubsection{Phase 3 (Long-term: verification)}
\label{sec:org03776c9}

\begin{center}
\begin{tabular}{lll}
Key & Type & Semantics\\
\hline
\texttt{:refines} & \texttt{datum} (result -> Bool) & Refinement predicate, compiles to Sigma type\\
\texttt{:measure} & \texttt{datum} (args -> Nat) & Termination measure for recursion\\
\texttt{:decreases} & \texttt{symbol} or \texttt{datum} & What decreases on recursive calls\\
\texttt{:proof} & \texttt{datum} or \texttt{:auto} & Proof strategy hint for verification\\
\end{tabular}
\end{center}
\subsection{Mapping Spec Keys to Dependent Types}
\label{sec:org5069bc2}

This is the deep theoretical connection between spec metadata and Prologos's
type system. Every metadata key has a \emph{type-theoretic interpretation}.
\subsubsection{Properties as Dependent Types}
\label{sec:orgf369ca8}

A property "for all x : Nat, add x 0 = x" is a Pi type:

\begin{verbatim}
;; The property:
:properties
  - :name "right-identity"
    :forall {x : Nat}
    :holds [eq? [add x 0N] x]

;; Corresponds to the dependent type:
;; <(x : Nat) -> [Eq Nat [add x 0N] x]>
;; A function from any Nat to a proof of equality
\end{verbatim}
\subsubsection{Preconditions as Refined Arguments}
\label{sec:org43938ab}

\begin{verbatim}
;; The contract:
spec divide Int Int -> Int
  :pre (fn [x y] [not [eq? y 0]])

;; Maps to the dependent type:
;; <(x : Int) -> (y : Int) -> [Not [Eq Int y 0]] -> Int>
;; Third argument is a proof that y /= 0
\end{verbatim}
\subsubsection{Postconditions as Refined Return Types}
\label{sec:org96f074c}

\begin{verbatim}
;; The contract:
spec abs Int -> Int
  :post (fn [x result] [>= result 0])

;; Maps to:
;; <(x : Int) -> <(result : Int) * [GTE result 0]>>
;; Return type is a Sigma pair: the result AND a proof it's non-negative
\end{verbatim}
\subsubsection{Invariants as Dependent Function Types}
\label{sec:org39e4f41}

\begin{verbatim}
;; Clojure spec's :fn relation:
spec ranged-rand Int Int -> Int
  :invariant (fn [start end result]
    [and [>= result start] [< result end]])

;; Maps to:
;; <(start : Int) -> (end : Int) -> {v : Int | v >= start, v < end}>
\end{verbatim}
\subsubsection{The Graduation Path}
\label{sec:orge6b6e58}

The key design insight: \emph{the same spec syntax upgrades from testing to proving
without modification}. At Phase 2, properties/contracts are runtime checks
(QuickCheck generators from types, contract wrappers at call boundaries). At
Phase 3, the same properties become proof obligations that the type checker
verifies statically. The syntax is identical; the verification backend upgrades.

\begin{center}
\begin{tabular}{rlll}
Phase & \texttt{:properties} & \texttt{:pre} / \texttt{:post} & \texttt{:refines}\\
\hline
2 & Runtime QuickCheck & Runtime contracts & Runtime assertion\\
3 & Proof obligation & Refined Pi type & Sigma compilation\\
\end{tabular}
\end{center}
\section{Part III: Typed Holes (\texttt{??})}
\label{sec:orga27a4ca}

\subsection{Motivation}
\label{sec:org8ee21f8}

Agda and Idris demonstrate that \emph{holes} are the foundation of interactive
development in dependently-typed languages. The programmer writes what they know,
marks what they don't with a hole, and the type system reports what's needed.

Prologos already has \texttt{\_} (inference hole / wildcard), which means "infer this
for me automatically." The \texttt{??} typed hole means something fundamentally
different: "I don't know what goes here --- show me what's possible."
\subsection{Syntax}
\label{sec:org9947416}

\begin{verbatim}
;; Unnamed hole:
defn reverse
  | [nil]        -> nil
  | [[cons h t]] -> ??

;; Named hole:
defn zip-with [f xs ys]
  match [xs ys]
    | [nil _]              -> nil
    | [_ nil]              -> nil
    | [[cons x xs'] [cons y ys']] -> [cons ??combine ??recurse]
\end{verbatim}
\subsubsection{Reader}
\label{sec:org2a8b72a}

The WS reader recognizes \texttt{??} (and \texttt{??identifier}) as a token type
\texttt{'typed-hole}. The form-builder emits \texttt{(\$typed-hole)} or \texttt{(\$typed-hole name)}
sentinel. The parser converts this to \texttt{expr-typed-hole} (new AST node).
\subsubsection{Distinction from \texttt{\_}}
\label{sec:orgf283c31}

\begin{center}
\begin{tabular}{lll}
Feature & \texttt{\_} (inference hole) & \texttt{??} (typed hole)\\
\hline
Intent & "Infer this for me" & "Help me fill this in"\\
Behavior & Silent, automatic & Interactive, diagnostic\\
Output & Resolved type/value & Hole report to editor\\
Use case & Type inference & Development workflow\\
\end{tabular}
\end{center}
\subsubsection{Hole Report}
\label{sec:orge5242be}

When the type checker encounters \texttt{expr-typed-hole}, it produces a diagnostic
(not an error) containing:

\begin{verbatim}
Hole ??combine : [List A]
Context:
  x  : A
  xs' : [List A]
  y  : B
  ys' : [List B]
  f  : A -> B -> C
  zip-with : [A -> B -> C] -> [List A] -> [List B] -> [List C]
Suggestions:
  [f x y]            -- from applying f to available arguments
  [cons [f x y] ...]  -- from matching the expected List type
\end{verbatim}

The editor protocol (future work) would transmit this as structured data.
\subsection{"The What?What?" Interaction Model}
\label{sec:org01d7ed4}

The envisioned workflow:

\begin{enumerate}
\item Write \texttt{spec} with type signature, properties, examples
\item Write \texttt{defn} skeleton with \texttt{??} holes for unknown parts
\item Press key command --- the type system reports what each hole needs
\item Select a suggestion or refine manually
\item Repeat until all holes are filled and properties hold
\end{enumerate}

This is fundamentally the Agda/Idris interaction model, but using Prologos's
structured spec metadata to provide richer suggestions (e.g., if a property says
the function is commutative, the case-split suggestions can leverage that).
\section{Part IV: The \texttt{property} Keyword}
\label{sec:org0e1a952}

\subsection{Motivation}
\label{sec:org4772d74}

Many functions share the same properties. Sorting functions are idempotent.
Monoids have identity and associativity. Functors preserve composition. Writing
these properties inline on every \texttt{spec} is repetitive and error-prone.

The \texttt{property} keyword provides reusable, composable proposition groups ---
analogous to \texttt{bundle} for traits. The parallel is exact:

\begin{center}
\begin{tabular}{lll}
Concept & For methods & For propositions\\
\hline
Single declaration & \texttt{trait} & individual \texttt{:holds} clause\\
Named group & \texttt{bundle} & \texttt{property}\\
Composition & \texttt{bundle} includes traits/bundles & \texttt{property} includes properties\\
Attachment to spec & \texttt{:where (Ord A)} & \texttt{:properties (sortable-laws A)}\\
Attachment to trait & (implicit) & \texttt{:laws (functor-laws F)}\\
\end{tabular}
\end{center}
\subsection{Syntax}
\label{sec:org81e4d5a}

\subsubsection{Basic Declaration}
\label{sec:org1de4fc8}

\begin{verbatim}
property sortable-laws {A : Type}
  :where (Ord A)
  - :name "idempotent"
    :forall {xs : [List A]}
    :holds [eq? [sort [sort xs]] [sort xs]]
  - :name "length-preserving"
    :forall {xs : [List A]}
    :holds [eq? [length [sort xs]] [length xs]]
  - :name "ordered-output"
    :forall {xs : [List A]}
    :holds [sorted? [sort xs]]
\end{verbatim}

A \texttt{property} is a named, parameterized conjunction of propositions. Each clause
is a universally quantified boolean assertion. The \texttt{:where} constrains what
types the property is meaningful for.
\subsubsection{Usage in \texttt{spec}}
\label{sec:orgfc98e01}

\begin{verbatim}
spec sort [List A] -> [List A]
  :implicits {A : Type}
  :where (Ord A)
  :properties (sortable-laws A)
\end{verbatim}

The \texttt{(sortable-laws A)} reference is like \texttt{(Ord A)} in a \texttt{:where} --- a named
requirement, parameterized by type variables from the spec.
\subsection{Composition via \texttt{:includes}}
\label{sec:org95991e5}

Properties compose via conjunction, paralleling \texttt{bundle}:

\begin{verbatim}
property semigroup-laws {A : Type}
  :where (Add A)
  - :name "associativity"
    :forall {x y z : A}
    :holds [eq? [add [add x y] z] [add x [add y z]]]

property monoid-laws {A : Type}
  :where (Add A) (AdditiveIdentity A)
  :includes (semigroup-laws A)
  - :name "left-identity"
    :forall {x : A}
    :holds [eq? [add additive-identity x] x]
  - :name "right-identity"
    :forall {x : A}
    :holds [eq? [add x additive-identity] x]

property group-laws {A : Type}
  :where (Add A) (AdditiveIdentity A) (Neg A)
  :includes (monoid-laws A)
  - :name "left-inverse"
    :forall {x : A}
    :holds [eq? [add [neg x] x] additive-identity]
  - :name "right-inverse"
    :forall {x : A}
    :holds [eq? [add x [neg x]] additive-identity]
\end{verbatim}

\texttt{group-laws} includes \texttt{monoid-laws} which includes \texttt{semigroup-laws}. Flattened:
\texttt{group-laws} is the conjunction of all six properties.
\subsection{The Logic Programming Reading}
\label{sec:org3c61716}

From the sequent calculus perspective:

\begin{verbatim}
property P {A} :where (C₁ A) (C₂ A) :includes (Q A) =
  Γ, C₁(A), C₂(A) ⊢ Q(A) ∧ p₁(A) ∧ p₂(A) ∧ ...
\end{verbatim}

Each \texttt{:holds} clause is a proposition. \texttt{:includes} is conjunction. The whole
\texttt{property} is a universally quantified sequent: "for all A satisfying C₁ and C₂,
the conjunction of all these propositions holds."

When attached to a spec with \texttt{:properties (P A)}, you assert that your
implementation satisfies that sequent. At Phase 2 (QuickCheck), each conjunct
becomes a test. At Phase 3 (verification), each conjunct becomes a proof
obligation.

The composability is clean because conjunction is associative and commutative ---
the order and nesting of \texttt{:includes} don't matter; the result is always a flat
set of propositions.
\subsection{Parameterized Properties (Higher-Order)}
\label{sec:org04619f5}

Properties can be parameterized over functions, not just types:

\begin{verbatim}
;; Properties about ANY sorting function
property sorting-properties {A : Type} {f : [List A] -> [List A]}
  :where (Ord A)
  - :name "idempotent"
    :forall {xs : [List A]}
    :holds [eq? [f [f xs]] [f xs]]
  - :name "permutation"
    :forall {xs : [List A]}
    :holds [same-elements? xs [f xs]]
  - :name "ordered"
    :forall {xs : [List A]}
    :holds [sorted? [f xs]]

;; Attach to multiple implementations:
spec sort [List A] -> [List A]
  :implicits {A : Type}
  :where (Ord A)
  :properties (sorting-properties A sort)

spec merge-sort [List A] -> [List A]
  :implicits {A : Type}
  :where (Ord A)
  :properties (sorting-properties A merge-sort)
\end{verbatim}

Define the contract once, attach it to any implementation.
\subsection{\texttt{:laws} on Traits}
\label{sec:org2a67970}

A trait can reference its laws:

\begin{verbatim}
property functor-laws {F : Type -> Type}
  :where (Functor F)
  - :name "identity"
    :forall {xs : [F A]}
    :holds [eq? [fmap id xs] xs]
  - :name "composition"
    :forall {f : [A -> B]} {g : [B -> C]} {xs : [F A]}
    :holds [eq? [fmap [>> f g] xs] [fmap g [fmap f xs]]]

trait Functor {F : Type -> Type}
  :laws (functor-laws F)
  spec fmap {A B : Type} [A -> B] -> [F A] -> [F B]
\end{verbatim}

\texttt{:laws} says "any instance of this trait is expected to satisfy these
properties." When you write \texttt{instance Functor List}, the system knows to check
\texttt{(functor-laws List)}. At Phase 2: generate \texttt{List}-specialized QuickCheck tests.
At Phase 3: proof obligations.
\subsection{Hierarchical Laws (Functor → Applicative → Monad)}
\label{sec:org56b910a}

\begin{verbatim}
property applicative-laws {F : Type -> Type}
  :where (Applicative F)
  :includes (functor-laws F)
  - :name "identity"
    :forall {xs : [F A]}
    :holds [eq? [ap [pure id] xs] xs]
  - :name "homomorphism"
    :forall {f : [A -> B]} {x : A}
    :holds [eq? [ap [pure f] [pure x]] [pure [f x]]]

property monad-laws {M : Type -> Type}
  :where (Monad M)
  :includes (applicative-laws M)
  - :name "left-identity"
    :forall {a : A} {f : [A -> [M B]]}
    :holds [eq? [bind [pure a] f] [f a]]
  - :name "right-identity"
    :forall {ma : [M A]}
    :holds [eq? [bind ma pure] ma]
  - :name "associativity"
    :forall {ma : [M A]} {f : [A -> [M B]]} {g : [B -> [M C]]}
    :holds [eq? [bind [bind ma f] g] [bind ma [fn [a] [bind [f a] g]]]]
\end{verbatim}
\subsection{Name Scoping in \texttt{:includes}}
\label{sec:orgef070df}

When \texttt{monad-laws} includes \texttt{applicative-laws} which includes \texttt{functor-laws},
and both \texttt{monad-laws} and \texttt{functor-laws} define a property named \texttt{"identity"}:
names are scoped to their declaring \texttt{property} block.

Flattened property names use \texttt{/} qualification:
\texttt{functor-laws/identity}, \texttt{applicative-laws/identity}, \texttt{monad-laws/left-identity}.

This parallels how trait methods are scoped to their trait.
\subsection{The Symmetry}
\label{sec:org7b72b98}

\begin{verbatim}
trait     : Type → Set of method signatures     (what you must implement)
bundle    : conjunction of traits                (shorthand for multiple traits)
property  : Type → Set of propositions           (what must be true)
  includes: conjunction of properties            (shorthand for multiple properties)

spec      : function → type + metadata
  :where  : trait constraints                    (what the types must support)
  :properties : property constraints             (what the function must satisfy)

instance  : Type → method implementations        (how you implement)
  (auto)  : Type → property evidence             (how you prove / test)
\end{verbatim}

Traits say what operations exist. Properties say what laws they obey. Today only
the first is checked; this design gives us a path to checking the second.
\section{Part V: Expressing Higher-Kinded Pi Types Ergonomically}
\label{sec:orgeeb260e}

\subsection{The Problem}
\label{sec:orgae55810}

Complex higher-kinded and higher-rank types are unreadable:

\begin{verbatim}
;; Current xf-compose signature:
spec xf-compose {A B C : Type} <(S :0 Type) -> [S -> B -> S] -> S -> A -> S> <(S :0 Type) -> [S -> C -> S] -> S -> B -> S> -> <(S :0 Type) -> [S -> C -> S] -> S -> A -> S>
\end{verbatim}

The transducer module already identifies the abstraction in a comment:
\texttt{Type: Xf A B = forall R. (R -> B -> R) -> (R -> A -> R)}. With the right name,
the signature becomes: \texttt{[Xf A B] -> [Xf B C] -> [Xf A C]} --- category-theoretic
composition, instantly readable.
\subsection{Progressive Complexity Levels}
\label{sec:org0cad009}

\subsubsection{Level 0: No types (inference only)}
\label{sec:orgf108ba9}

\begin{verbatim}
defn add [x y] [nat-add x y]
\end{verbatim}
\subsubsection{Level 1: Simple spec (monomorphic)}
\label{sec:org1f36fc0}

\begin{verbatim}
spec add Nat Nat -> Nat
defn add [x y] [nat-add x y]
\end{verbatim}
\subsubsection{Level 2: Polymorphic spec}
\label{sec:orgb11e810}

\begin{verbatim}
spec id A -> A
defn id [x] x
\end{verbatim}

(\texttt{\{A : Type\}} inferred from \texttt{A} appearing free in signature)
\subsubsection{Level 3: Constrained polymorphism}
\label{sec:orgdaba1b2}

\begin{verbatim}
spec sort [List A] -> [List A]
  :where (Ord A)
defn sort [xs] ...
\end{verbatim}
\subsubsection{Level 4: Higher-kinded polymorphism}
\label{sec:org38b8e50}

\begin{verbatim}
spec fmap [A -> B] -> [F A] -> [F B]
  :where (Functor F)
  :doc "Apply a function inside a functor"
  :properties (functor-laws F)
defn fmap [f xs] ...
\end{verbatim}

(\texttt{\{A B : Type\}} inferred; \texttt{\{F : Type -> Type\}} inferred from \texttt{(Functor F)}
where clause --- the trait declares the kind)
\subsubsection{Level 5: Named type abstractions (functor)}
\label{sec:org9f24479}

\begin{verbatim}
spec xf-compose [Xf A B] -> [Xf B C] -> [Xf A C]
  :implicits {A B C : Type}
  :properties (transducer-fusion-laws A B C)
\end{verbatim}

(\texttt{Xf} defined via \texttt{functor} --- see Part VII)
\subsubsection{Level 6: Dependent types (Pi)}
\label{sec:org037b864}

\begin{verbatim}
spec replicate <(n : Nat) -> A -> [Vec A n]>
  :implicits {A : Type}
  :doc "Create a vector of n copies of a value"
  :examples
    - [replicate 3N :x] => @[:x :x :x]
  :properties
    - :name "correct-length"
      :forall {n : Nat} {a : A}
      :holds [eq? [length [replicate n a]] n]
\end{verbatim}
\subsubsection{Level 7: Dependent types (Sigma / existential)}
\label{sec:org7e6c38f}

\begin{verbatim}
spec filter [A -> Bool] -> [List A] -> <(result : [List A]) * [SubList result xs]>
  :implicits {A : Type}
  :doc "Filter elements, with proof that result is a sublist"
\end{verbatim}
\subsection{Key Ergonomic Decisions}
\label{sec:orgcdf26c6}

\begin{enumerate}
\item \textbf{Pi/Sigma types never need to leak to the surface}: \texttt{functor} declarations
name complex type abstractions with approachable metadata. The underlying
Pi type lives in \texttt{:unfolds}, consulted only when needed.

\item \textbf{Angle brackets contain the complexity}: \texttt{<...>} delimits the dependent part
of a type. Everything outside angle brackets is familiar ML-style syntax.

\item \textbf{Metadata explains the theory}: The \texttt{:doc} key provides human-readable
explanation. Properties double as documentation and specification.

\item \textbf{Examples ground the abstraction}: Every complex type should have concrete
examples. The \texttt{:examples} key makes this a first-class concern.

\item \textbf{Where-clauses read as English}: "sort takes a list of A where A has ordering"
is near-English. The trait constraint system hides the dictionary-passing
mechanism.

\item \textbf{Properties express what the function \emph{means}}: Rather than encoding
everything in the type (which can be opaque), properties express semantic
invariants in a universally readable format.

\item \textbf{\texttt{:implicits} declutters signatures}: Implicit binders move to metadata when
they would add visual weight to the type signature.
\end{enumerate}
\section{Part VI: Improved Implicit Inference}
\label{sec:org1651ace}

Today, \texttt{\{A : Type\}} is required even when \texttt{A} appears in the type signature
(e.g., \texttt{[List A] -> Nat}). The type checker can already infer this in many
cases. Two improvements are proposed:
\subsection{Direction 1: Auto-introduce unbound type variables}
\label{sec:orgc330027}

If a capitalized identifier \texttt{A} appears free in the type signature and is not
bound by any explicit binder, auto-introduce \texttt{\{A : Type\}}. This already works
for simple cases; the proposal is to make it the documented, reliable behavior.

\begin{verbatim}
;; These would be equivalent:
spec length {A : Type} [List A] -> Nat
spec length [List A] -> Nat
\end{verbatim}
\subsection{Direction 2: Kind inference from \texttt{:where} clauses}
\label{sec:org01efac9}

If \texttt{C} appears in \texttt{:where (Seqable C)} and \texttt{Seqable} is declared over
\texttt{\{C : Type -> Type\}}, infer \texttt{C}'s kind from the where clause. This eliminates
the most annoying case --- having to write \texttt{\{C : Type -> Type\}} just because
you use an HKT trait.

\begin{verbatim}
;; These would be equivalent:
spec gmap {A B : Type} {C : Type -> Type} [A -> B] -> [C A] -> [C B]
  :where (Seqable C) (Buildable C)

spec gmap [A -> B] -> [C A] -> [C B]
  :where (Seqable C) (Buildable C)
\end{verbatim}

Note: \texttt{propagate-kinds-from-constraints} already does this internally. The
proposal is to make it work when explicit \texttt{\{...\}} binders are omitted entirely.
\subsection{Explicit Always Available}
\label{sec:orge509b73}

Explicit binders remain for: pedagogic clarity, disambiguation when inference
produces unexpected results, and the rare cases where no constraining position
exists (e.g., \texttt{spec empty \{A : Type\} [List A]}).
\section{Part VII: The \texttt{functor} Keyword}
\label{sec:org643c950}

\subsection{Motivation}
\label{sec:org8ecdc2a}

The \texttt{functor} keyword provides named type abstractions with category-theoretic
metadata. It solves the fundamental readability problem of higher-kinded and
higher-rank types by giving them \emph{meaning through structure}, not just through
their expansion.

The design philosophy: you understand \texttt{Xf A B} by knowing what it \emph{does} (":doc
A transducer"), how it \emph{combines} (":compose xf-compose"), and what \emph{rules} it
follows (":laws transducer-fusion-laws"). The Pi type it unfolds to is available
but is consulted last, not first.
\subsection{Syntax}
\label{sec:org58e0864}

\begin{verbatim}
functor Xf {A B : Type}
  :doc "A transducer: transforms A-reductions into B-reductions"
  :compose xf-compose
  :identity id-xf
  :laws (transducer-fusion-laws A B)
  :unfolds <(S :0 Type) -> [S -> B -> S] -> S -> A -> S>
\end{verbatim}
\subsection{Progressive Disclosure}
\label{sec:org96ac7a6}

Like \texttt{spec}, the metadata keys are progressive. Simple synonyms are minimal:

\begin{verbatim}
functor FilePath
  :unfolds String
\end{verbatim}

Parameterized synonyms add type parameters:

\begin{verbatim}
functor Result {A}
  :unfolds [Either String A]
  :doc "A computation that may fail with a string error"
\end{verbatim}

Algebraically structured types add operations:

\begin{verbatim}
functor Xf {A B : Type}
  :doc "A transducer: transforms A-reductions into B-reductions"
  :compose xf-compose
  :identity id-xf
  :unfolds <(S :0 Type) -> [S -> B -> S] -> S -> A -> S>
\end{verbatim}

Full category-theoretic treatment adds laws:

\begin{verbatim}
functor Lens {S T A B : Type}
  :doc "A bidirectional accessor: view and update a part of a structure"
  :compose lens-compose
  :identity lens-id
  :laws (lens-laws S T A B)
  :see-also [Prism Traversal Iso]
  :unfolds <{F : Type -> Type} -> (Functor F) -> [A -> [F B]] -> S -> [F T]>
\end{verbatim}
\subsection{Category-Theoretic Metadata Keys}
\label{sec:orgf638b8b}

The keys map CT concepts to approachable language:

\begin{center}
\begin{tabular}{lll}
CT concept & Key & Meaning in plain English\\
\hline
Object mapping & \texttt{:unfolds} & "What this actually expands to under the hood"\\
Morphism composition & \texttt{:compose} & "How to chain two of these together"\\
Identity morphism & \texttt{:identity} & "The do-nothing version"\\
Laws & \texttt{:laws} & "What rules these always follow"\\
Documentation & \texttt{:doc} & "What this means in human terms"\\
Natural transformation & \texttt{:transforms} & "How to convert between different functors" (future)\\
Adjunction & \texttt{:adjoint} & "What functor is this paired with" (far future)\\
\end{tabular}
\end{center}

Only \texttt{:unfolds} is required (for Phase 1). The rest are progressive. A user
who never thinks about category theory writes:

\begin{verbatim}
functor Xf {A B : Type}
  :doc "A transducer"
  :unfolds <(S :0 Type) -> [S -> B -> S] -> S -> A -> S>
\end{verbatim}

A user who cares about algebraic structure adds \texttt{:compose} and \texttt{:identity}. A
user who wants verification adds \texttt{:laws}. The CT is always there but only
surfaces when you reach for it.
\subsection{Why \texttt{functor} and Not \texttt{abbrev} or \texttt{type}}
\label{sec:org026124d}

\texttt{abbrev} doesn't capture the ambition. An abbreviation is just a shorter name
for the same thing. What we're proposing is a \emph{named abstraction with declared
structure} --- the operations, laws, documentation are first-class. \texttt{abbrev}
makes you think of \texttt{typedef} --- mechanical shortening. \texttt{functor} makes you
think of structure-preserving mappings --- which is what these actually are.

\texttt{type} clashes with \texttt{deftype} (which creates algebraic data types with
constructors). A \texttt{functor} doesn't create constructors; it names a type with
structure.

The \texttt{functor} (keyword) vs \texttt{Functor} (trait) naming is consistent with Prologos's
existing pattern: all keywords are lowercase (\texttt{spec}, \texttt{defn}, \texttt{trait}, \texttt{bundle});
all type/trait names are capitalized. The relationship: a \texttt{functor} declaration
might \emph{itself} be a \texttt{Functor} instance (in the trait sense). The keyword names
the concept; the trait constrains it.
\subsection{Transparency Model}
\label{sec:org33b0c02}

\texttt{functor} declarations are transparent by default (like Lean's \texttt{abbrev}). The
type checker unfolds \texttt{Xf A B} to its \texttt{:unfolds} form during elaboration. Trait
instances carry through: if \texttt{Eq (Xf Nat Bool)} is needed and \texttt{Eq} has an
instance for the expanded type, it resolves.

Future consideration: opaque functors (no \texttt{:unfolds}, only described by
operations) for abstract type boundaries. Deferred to post-Phase 1.
\subsection{Error Messages}
\label{sec:orgcad4477}

Error messages show the \texttt{functor} name by default:

\begin{verbatim}
Error: Expected [Xf Nat Bool], got [Xf Nat Nat]
\end{verbatim}

Not the expanded form. An \texttt{-{}-{}expand-types} flag or editor hover reveals the
full type when needed.
\subsection{How \texttt{functor} Composes with \texttt{spec}, \texttt{property}, and \texttt{trait}}
\label{sec:org8323b77}

\begin{verbatim}
;; The type abstraction
functor Xf {A B : Type}
  :doc "A transducer: transforms A-reductions into B-reductions"
  :compose xf-compose
  :identity id-xf
  :laws (transducer-fusion-laws A B)
  :unfolds <(S :0 Type) -> [S -> B -> S] -> S -> A -> S>

;; Reusable properties (using the functor name!)
property transducer-fusion-laws {A B C : Type}
  - :name "compose-associativity"
    :forall {f : [Xf A B]} {g : [Xf B C]} {h : [Xf C D]}
    :holds [eq? [xf-compose [xf-compose f g] h] [xf-compose f [xf-compose g h]]]
  - :name "left-identity"
    :forall {f : [Xf A B]}
    :holds [eq? [xf-compose id-xf f] f]
  - :name "right-identity"
    :forall {f : [Xf A B]}
    :holds [eq? [xf-compose f id-xf] f]

;; Function specs using the functor name
spec xf-compose [Xf A B] -> [Xf B C] -> [Xf A C]
  :implicits {A B C : Type}
  :doc "Compose two transducers (apply first argument first)"
  :properties (transducer-fusion-laws A B C)

spec map-xf [A -> B] -> [Xf A B]
  :implicits {A B : Type}
  :doc "Transform each element through f before passing to reducer"

spec filter-xf [A -> Bool] -> [Xf A A]
  :implicits {A : Type}
  :doc "Only pass elements satisfying pred to the reducer"

spec transduce [Xf A B] -> [R -> B -> R] -> R -> [List A] -> R
  :implicits {A B R : Type}
  :doc "Apply a transducer to a list with a reducer and initial accumulator"

spec into-list [Xf A B] -> [List A] -> [List B]
  :implicits {A B : Type}
  :doc "Transduce into a correctly-ordered list"
  :examples
    - [into-list [map-xf suc] '[1N 2N 3N]] => '[2N 3N 4N]
\end{verbatim}

The \texttt{functor} names the type. The \texttt{property} states its laws. The \texttt{spec} uses
both. Everything is readable. The Pi types are present exactly once, in
\texttt{:unfolds}, consulted only when necessary.
\section{Part VIII: Implementation Considerations}
\label{sec:orgf1312d7}

\subsection{spec-entry Struct Extension}
\label{sec:org319a199}

The recommended approach: add a single \texttt{metadata} field that stores the entire
keyword map as a hash table. Individual accessors extract from the metadata hash.
This avoids struct migration pain:

\begin{verbatim}
(struct spec-entry
  (type-datums docstring multi? srcloc where-constraints
   implicit-binders rest-type
   metadata)  ;; <-- one new field: hasheq of keyword -> datum
  #:transparent)

;; Accessor helpers:
(define (spec-entry-examples e)
  (hash-ref (or (spec-entry-metadata e) (hasheq)) ':examples '()))
(define (spec-entry-properties e)
  (hash-ref (or (spec-entry-metadata e) (hasheq)) ':properties '()))
(define (spec-entry-implicits-meta e)
  (hash-ref (or (spec-entry-metadata e) (hasheq)) ':implicits #f))
\end{verbatim}
\subsection{rewrite-implicit-map Extension}
\label{sec:orgd4ef215}

Currently scoped to \texttt{def} and \texttt{defn} heads (macros.rkt line 2452). Extending to
\texttt{spec} is a one-line change in the \texttt{memq} check. However, the spec head requires
special handling: the type signature tokens must be separated from the trailing
keyword map. This can be done by scanning for the first keyword-headed child
whose indentation matches the spec's indentation level.
\subsection{Typed Hole AST Node}
\label{sec:org4961345}

\begin{verbatim}
(struct expr-typed-hole (name) #:transparent)  ;; name is #f or symbol
\end{verbatim}

Type checker (typing-core.rkt): when encountering \texttt{expr-typed-hole}, compute
the expected type from the context and emit a diagnostic report rather than
returning \texttt{(expr-error)}.
\subsection{property-entry Store}
\label{sec:org40de4d9}

\begin{verbatim}
(struct property-entry
  (name           ;; symbol
   params         ;; alist of ((name . kind) ...)
   where-clauses  ;; (listof constraint)
   includes       ;; (listof (property-ref ...))
   clauses        ;; (listof property-clause)
   metadata)      ;; hasheq for extensibility
  #:transparent)

(struct property-clause
  (name           ;; string (human-readable label)
   forall-binders ;; alist of ((name . type) ...)
   holds-expr)    ;; datum (boolean expression)
  #:transparent)

(define current-property-store (make-parameter (hasheq)))
\end{verbatim}

\texttt{:includes} references are resolved at registration time by looking up the
included property in the store and appending its (flattened) clauses. Clause
names are prefixed with the declaring property's name and \texttt{/}.
\subsection{functor-entry Store}
\label{sec:orgc69b8df}

\begin{verbatim}
(struct functor-entry
  (name         ;; symbol (capitalized)
   params       ;; alist of ((name . kind) ...)
   unfolds      ;; type datum (the expansion)
   metadata)    ;; hasheq (:compose, :identity, :laws, :doc, etc.)
  #:transparent)

(define current-functor-store (make-parameter (hasheq)))
\end{verbatim}

During elaboration, when the elaborator encounters a capitalized identifier
followed by arguments that matches a registered functor, it expands to the
\texttt{unfolds} form with parameters substituted.
\subsection{\texttt{:implicits} Merging in process-spec}
\label{sec:orgdc7080a}

\begin{verbatim}
;; In process-spec, after extracting metadata:
(define meta-implicits
  (if metadata
      (extract-binders-from-implicits-key metadata)
      '()))
(define merged-implicits
  (deduplicate-binders (append inline-implicits meta-implicits)))
;; Warn on duplicates between inline and :implicits
\end{verbatim}
\subsection{Grammar Extensions}
\label{sec:org93369ee}

Additions to \texttt{grammar.ebnf}:

\begin{verbatim}
(* Typed holes for interactive development *)
typed-hole      = '??' , [ identifier ] ;

(* Extended spec with optional metadata *)
spec-form       = 'spec' , identifier , spec-body , [ spec-metadata ] ;
spec-metadata   = spec-meta-entry , { spec-meta-entry } ;
spec-meta-entry = keyword-literal , spec-meta-value ;
spec-meta-value = expr
                | dash-list
                | spec-meta-map
                ;

(* Property declarations *)
property-form   = 'property' , identifier , [ implicit-binders ]
                , [ property-metadata ] , { property-clause } ;
property-metadata = property-meta-entry , { property-meta-entry } ;
property-meta-entry = ':where' , '(' , constraint , { constraint } , ')'
                    | ':includes' , '(' , property-ref , { property-ref } , ')'
                    ;
property-clause = '-' , ':name' , string-literal
                , ':forall' , '{' , binder , { binder } , '}'
                , ':holds' , expr ;

(* Functor declarations *)
functor-form    = 'functor' , identifier , [ implicit-binders ]
                , functor-metadata ;
functor-metadata = functor-meta-entry , { functor-meta-entry } ;
functor-meta-entry = ':unfolds' , type-expr
                   | ':compose' , identifier
                   | ':identity' , identifier
                   | ':laws' , '(' , property-ref , { property-ref } , ')'
                   | ':doc' , string-literal
                   | ':see-also' , '[' , identifier , { identifier } , ']'
                   ;
\end{verbatim}
\section{Part IX: Worked Examples}
\label{sec:orge0fe50e}

\subsection{Example 1: Simple spec with documentation, examples, and reusable properties}
\label{sec:org6554574}

\begin{verbatim}
property cons-head-law {A : Type}
  - :name "head-of-cons"
    :forall {x : A} {xs : [List A]}
    :holds [eq? [head [cons x xs]] [just x]]

spec head [List A] -> A?
  :implicits {A : Type}
  :doc "Returns the first element, or nothing if empty"
  :examples
    - [head '[1N 2N 3N]] => [just 1N]
    - [head '[]] => nothing
  :properties (cons-head-law A)

defn head
  | [nil]        -> nothing
  | [[cons h _]] -> [just h]
\end{verbatim}
\subsection{Example 2: Contract-style spec}
\label{sec:orgd624775}

\begin{verbatim}
spec divide Int Int -> Int
  :pre (fn [_ y] [not [eq? y 0]])
  :post (fn [x y r] [eq? [+ [* r y] [int-mod x y]] x])
  :doc "Integer division; undefined for zero divisor"

defn divide [x y] [int-div x y]
\end{verbatim}
\subsection{Example 3: Higher-kinded functor laws (using property + :laws)}
\label{sec:org4402f63}

\begin{verbatim}
property functor-laws {F : Type -> Type}
  :where (Functor F)
  - :name "identity-law"
    :forall {xs : [F A]}
    :holds [eq? [fmap id xs] xs]
  - :name "composition-law"
    :forall {f : [A -> B]} {g : [B -> C]} {xs : [F A]}
    :holds [eq? [fmap [>> f g] xs] [fmap g [fmap f xs]]]

trait Functor {F : Type -> Type}
  :laws (functor-laws F)
  spec fmap {A B : Type} [A -> B] -> [F A] -> [F B]

spec fmap [A -> B] -> [F A] -> [F B]
  :where (Functor F)
  :properties (functor-laws F)
\end{verbatim}
\subsection{Example 4: Transducer module (functor + property + spec)}
\label{sec:orge98cc3f}

\begin{verbatim}
ns prologos::data::transducer

functor Xf {A B : Type}
  :doc "A transducer: transforms A-reductions into B-reductions"
  :compose xf-compose
  :identity id-xf
  :laws (transducer-fusion-laws A B)
  :unfolds <(S :0 Type) -> [S -> B -> S] -> S -> A -> S>

property transducer-fusion-laws {A B C : Type}
  - :name "compose-associativity"
    :forall {f : [Xf A B]} {g : [Xf B C]} {h : [Xf C D]}
    :holds [eq? [xf-compose [xf-compose f g] h] [xf-compose f [xf-compose g h]]]
  - :name "left-identity"
    :forall {f : [Xf A B]}
    :holds [eq? [xf-compose id-xf f] f]
  - :name "right-identity"
    :forall {f : [Xf A B]}
    :holds [eq? [xf-compose f id-xf] f]

spec map-xf [A -> B] -> [Xf A B]
  :implicits {A B : Type}
  :doc "Transform each element through f before passing to reducer"

spec filter-xf [A -> Bool] -> [Xf A A]
  :implicits {A : Type}
  :doc "Only pass elements satisfying pred to the reducer"

spec xf-compose [Xf A B] -> [Xf B C] -> [Xf A C]
  :implicits {A B C : Type}
  :doc "Compose two transducers (apply first argument first)"
  :properties (transducer-fusion-laws A B C)

spec transduce [Xf A B] -> [R -> B -> R] -> R -> [List A] -> R
  :implicits {A B R : Type}
  :doc "Apply a transducer to a list with a reducer and initial accumulator"

spec into-list [Xf A B] -> [List A] -> [List B]
  :implicits {A B : Type}
  :doc "Transduce into a correctly-ordered list"
  :examples
    - [into-list [map-xf suc] '[1N 2N 3N]] => '[2N 3N 4N]
\end{verbatim}
\subsection{Example 5: Interactive development with ??}
\label{sec:orge80ea14}

\begin{verbatim}
spec reverse [List A] -> [List A]
  :implicits {A : Type}
  :properties
    - :name "involution"
      :forall {xs : [List Nat]}
      :holds [eq? [reverse [reverse xs]] xs]
    - :name "length-preserving"
      :forall {xs : [List Nat]}
      :holds [eq? [length [reverse xs]] [length xs]]

defn reverse
  | [nil]        -> nil
  | [[cons h t]] -> ??
  ;; Type system reports:
  ;;   Hole ?? : [List A]
  ;;   Context: h : A, t : [List A], reverse : [List A] -> [List A]
  ;;   Suggestions: [append [reverse t] [cons h nil]]
\end{verbatim}
\subsection{Example 6: Algebraic property hierarchy}
\label{sec:org245c983}

\begin{verbatim}
property semigroup-laws {A : Type}
  :where (Add A)
  - :name "associativity"
    :forall {x y z : A}
    :holds [eq? [add [add x y] z] [add x [add y z]]]

property monoid-laws {A : Type}
  :where (Add A) (AdditiveIdentity A)
  :includes (semigroup-laws A)
  - :name "left-identity"
    :forall {x : A}
    :holds [eq? [add additive-identity x] x]
  - :name "right-identity"
    :forall {x : A}
    :holds [eq? [add x additive-identity] x]

;; Usage:
spec fold [A -> A -> A] -> A -> [List A] -> A
  :implicits {A : Type}
  :where (Add A) (AdditiveIdentity A)
  :properties (monoid-laws A)
\end{verbatim}
\subsection{Example 7: Optics (functor for Lens)}
\label{sec:org7d8861f}

\begin{verbatim}
functor Lens {S T A B : Type}
  :doc "A bidirectional accessor: view and update a part of a structure"
  :compose lens-compose
  :identity lens-id
  :laws (lens-laws S T A B)
  :see-also [Prism Traversal Iso]
  :unfolds <{F : Type -> Type} -> (Functor F) -> [A -> [F B]] -> S -> [F T]>

spec lens-view [Lens S S A A] -> S -> A
  :doc "Extract the focus from a structure"

spec lens-set [Lens S S A A] -> A -> S -> S
  :doc "Replace the focus in a structure"
  :properties
    - :name "set-get"
      :forall {l : [Lens S S A A]} {s : S} {a : A}
      :holds [eq? [lens-view l [lens-set l a s]] a]
\end{verbatim}
\subsection{Example 8: Refinement type (future)}
\label{sec:org2f6bdb2}

\begin{verbatim}
spec abs Int -> Int
  :refines (fn [r] [>= r 0])
  :doc "Absolute value --- result is always non-negative"

;; Internally, :refines compiles to return type:
;; <(r : Int) * [GTE r 0]>
;; The function returns a pair: the result AND a proof of non-negativity
\end{verbatim}
\subsection{Example 9: Dependent-type-heavy spec}
\label{sec:org0fe95e4}

\begin{verbatim}
spec zip-with-vec <(n : Nat) -> [A -> B -> C] -> [Vec A n] -> [Vec B n] -> [Vec C n]>
  :implicits {A B C : Type}
  :doc "Zip two vectors of the same length with a combining function"
  :examples
    - [zip-with-vec 3N + @[1 2 3] @[10 20 30]] => @[11 22 33]
  :properties
    - :name "preserves-length"
      :forall {n : Nat} {f : [Int -> Int -> Int]}
              {xs : [Vec Int n]} {ys : [Vec Int n]}
      :holds [eq? [vec-length [zip-with-vec n f xs ys]] n]
\end{verbatim}
\section{Part X: Phased Roadmap}
\label{sec:org69b00ad}

\subsection{Phase 1: Syntax and Storage (Near-term)}
\label{sec:org173d39b}

\textbf{In scope now}:
\begin{itemize}
\item Extend \texttt{rewrite-implicit-map} to trigger on \texttt{spec} heads
\item Extend \texttt{process-spec} to extract keyword entries into metadata hash
\item Add \texttt{metadata} field to \texttt{spec-entry}
\item \texttt{:implicits} key: extract from metadata, merge with inline binders
\item \texttt{:doc} replaces positional docstring; \texttt{:where} augments positional \texttt{where}
\item Store \texttt{:examples}, \texttt{:see-also} in metadata
\item \texttt{??} reader token and \texttt{expr-typed-hole} AST node (report-only)
\item \texttt{property} keyword: parser, store, \texttt{:includes} flattening, name-scoped clauses
\item \texttt{:laws} key on \texttt{trait} declarations
\item \texttt{functor} keyword: parser, store, transparent expansion during elaboration
\item Grammar updates (grammar.ebnf, grammar.org)
\end{itemize}
\subsection{Phase 1b: Improved Implicit Inference}
\label{sec:org84cbb25}

\begin{itemize}
\item Auto-introduce unbound type variables as \texttt{\{X : Type\}}
\item Kind inference from \texttt{:where} clauses (extend \texttt{propagate-kinds-from-constraints}
to work when no explicit binders are present)
\item Goal: \texttt{\{A : Type\}} almost never needed for kind-\texttt{Type} variables;
\texttt{\{C : Type -> Type\}} unnecessary when \texttt{:where} pins the kind
\end{itemize}
\subsection{Phase 2: Example and Property Checking (Medium-term)}
\label{sec:org42f4b34}

\begin{itemize}
\item Parse \texttt{:examples} entries, type-check them, run them as tests
\item \texttt{Gen} trait for type-directed random value generation
\item Parse \texttt{:properties}, generate random inputs from types, check \texttt{:holds}
\item Property checking for \texttt{:laws} on trait instances
\item Contract wrapping: \texttt{:pre\textasciitilde{}/}:post\textasciitilde{} generate runtime checks at function
boundaries (Racket-contract style blame tracking)
\end{itemize}
\subsection{Phase 3: Refinement and Verification (Long-term)}
\label{sec:orgb81e8ed}

\begin{itemize}
\item \texttt{:refines} compiles to Sigma types (return type is \texttt{<(r : T) * P(r)>})
\item \texttt{:properties} become compile-time proof obligations
\item Proof search integration: \texttt{:proof :auto} triggers the logic engine
\item \texttt{:measure} for termination checking
\item Opaque functors (no \texttt{:unfolds}) for abstract type boundaries
\end{itemize}
\subsection{Phase 4: Interactive Theorem Proving (Far future)}
\label{sec:org90b3558}

\begin{itemize}
\item Editor protocol for \texttt{??} hole interaction
\item Case splitting via type information (Agda \texttt{C-c C-c})
\item Proof search (Agda \texttt{C-c C-a}) via propagator-based logic engine
\item Refinement reflection (Liquid Haskell style)
\item \texttt{:transforms} and \texttt{:adjoint} keys on \texttt{functor} for natural transformations
\end{itemize}
\section{Part XI: Resolved Design Questions}
\label{sec:org34386d9}

\begin{enumerate}
\item \textbf{Should \texttt{:where} in the map replace or augment positional \texttt{where}?}
\textbf{Decision}: augment (union of both constraint sets). Positional form
remains for backward compatibility; map form is preferred going forward.

\item \textbf{Should examples use \texttt{=>} or \texttt{->} as the separator?}
\textbf{Decision}: \texttt{=>} to avoid confusion with the arrow in type signatures.

\item \textbf{How should multi-arity specs handle metadata?}
\textbf{Decision}: metadata is function-level (shared across all arities). Each
\texttt{|} branch defines its own type; the metadata applies to the function as a
whole. Properties' \texttt{:forall} types determine which branch is exercised.
Examples' argument types determine dispatch.

\item \textbf{Should \texttt{:properties} quantified variables shadow spec's implicit binders?}
\textbf{Decision}: yes, with a warning. The property is a self-contained
expression; it should be readable without consulting the outer spec.

\item \textbf{What happens to unrecognized keys?}
\textbf{Decision}: store silently (forward-compatibility). Tool-specific keys
(e.g., \texttt{:editor-hint}, \texttt{:benchmark}) can be added without language changes.

\item \textbf{Is \texttt{:unfolds} always required on \texttt{functor}?}
\textbf{Decision}: yes, for Phase 1. Opaque functors (no \texttt{:unfolds}) deferred to
Phase 3.

\item \textbf{Should \texttt{functor} handle non-CT uses (plain type synonyms)?}
\textbf{Decision}: yes --- \texttt{functor} is the single keyword for "I'm naming a type
concept." Progressive metadata: no \texttt{:compose} means no algebraic structure.
\end{enumerate}
\section{References}
\label{sec:org9941de5}

\subsection{Clojure Spec}
\label{sec:orgd4d199e}
\begin{itemize}
\item \href{https://clojure.org/guides/spec}{Clojure Spec Guide}
\item \href{https://clojuredocs.org/clojure.spec.alpha/fdef}{s/fdef Reference}
\item \href{https://clojure.org/about/spec}{Spec Rationale and Overview}
\end{itemize}
\subsection{Malli}
\label{sec:org0d6918d}
\begin{itemize}
\item \href{https://github.com/metosin/malli}{Malli GitHub}
\item \href{https://github.com/metosin/malli/blob/master/docs/function-schemas.md}{Malli Function Schemas}
\end{itemize}
\subsection{PropEr}
\label{sec:org219b434}
\begin{itemize}
\item \href{https://github.com/proper-testing/proper}{PropEr GitHub}
\item \href{https://proper-testing.github.io/tutorials/PropEr\_introduction\_to\_Property-Based\_Testing.html}{PropEr Tutorial}
\end{itemize}
\subsection{Agda}
\label{sec:orge075587}
\begin{itemize}
\item \href{https://agda.readthedocs.io/en/latest/tools/auto.html}{Agda Auto Proof Search}
\item \href{https://agda.readthedocs.io/en/latest/getting-started/a-taste-of-agda.html}{A Taste of Agda}
\end{itemize}
\subsection{Idris 2}
\label{sec:org6842a4c}
\begin{itemize}
\item \href{https://idris2.readthedocs.io/en/latest/tutorial/interactive.html}{Idris 2 Interactive Editing}
\end{itemize}
\subsection{Lean 4}
\label{sec:org56182c0}
\begin{itemize}
\item \href{https://lean-lang.org/theorem\_proving\_in\_lean4/type\_classes.html}{Lean 4 Type Classes}
\item \href{https://lean-lang.org/doc/reference/latest/Type-Classes/Instance-Synthesis/}{Lean 4 Instance Synthesis}
\end{itemize}
\subsection{Liquid Haskell}
\label{sec:orgbff9176}
\begin{itemize}
\item \href{https://nikivazou.github.io/lh-course/Lecture\_01\_RefinementTypes.html}{Refinement Types Tutorial}
\item \href{https://goto.ucsd.edu/\~nvazou/refinement\_types\_for\_haskell.pdf}{Refinement Types for Haskell (paper)}
\end{itemize}
\subsection{Haskell Type Synonyms}
\label{sec:orgb1a9154}
\begin{itemize}
\item \href{https://wiki.haskell.org/Type\_synonym}{GHC Type Synonyms}
\end{itemize}
\subsection{Koka}
\label{sec:org027bd7d}
\begin{itemize}
\item \href{https://koka-lang.github.io/koka/doc/book.html}{Koka Language Book}
\end{itemize}
\subsection{F*}
\label{sec:org66dd166}
\begin{itemize}
\item \href{https://github.com/FStarLang/FStar/wiki/Qualifiers-for-definitions-and-declarations}{F* Qualifier Wiki}
\end{itemize}
\subsection{1ML}
\label{sec:org364a6ef}
\begin{itemize}
\item \href{https://people.mpi-sws.org/\~rossberg/1ml/1ml-extended.pdf}{1ML (Rossberg) — Core and Modules United}
\end{itemize}
\subsection{Racket Contracts}
\label{sec:orgd6ded2b}
\begin{itemize}
\item \href{https://docs.racket-lang.org/guide/contracts.html}{Racket Contracts Guide}
\end{itemize}
\end{document}
